\documentclass[aspectratio=169,10pt,t]{beamer}
\usepackage{ac1}
\aulainfo{33}{Refatoração e Testes}{Legibilidade, organização, robustez e boas práticas}

\begin{document}
\capa

\begin{frame}{Objetivos da aula}
  \begin{itemize}
    \item \textbf{Refatorar}: melhorar a legibilidade, a organização e a eficiência \emph{sem mudar o comportamento}
    \item \textbf{Testar} casos extremos para garantir a robustez
    \item Aplicar boas práticas: nomes, coesão, modularização e separação de responsabilidades
    \item Receber feedback detalhado sobre as melhorias e ajustes finais
  \end{itemize}
  \begin{exampleblock}{Cap.~9 em ação}
    Preocupe-se com seu leitor; seja simples; seja desapegado; teste seu programa.
  \end{exampleblock}
\end{frame}

\begin{frame}{O que é refatorar?}
  \begin{columns}[T]
    \column{0.52\textwidth}
    \begin{itemize}
      \item Reorganizar o código \textbf{sem alterar o que ele faz}
      \item Como garantir isso? \alert{Com testes}: rode-os antes e depois
      \item Faça mudanças \emph{pequenas}, uma por vez
    \end{itemize}
    \column{0.45\textwidth}
    \begin{block}{Fluxo seguro}
      \begin{enumerate}\small
        \item Garanta que os testes passam
        \item Faça \emph{uma} melhoria
        \item Rode os testes de novo
        \item Repita
      \end{enumerate}
    \end{block}
  \end{columns}
\end{frame}

\begin{frame}{Sinais de que é hora de refatorar}
  \begin{columns}[T]
    \column{0.5\textwidth}
    \begin{itemize}
      \item Funções com \(>\) 30 linhas ou que fazem várias coisas
      \item Código \textbf{repetido} (copiar e colar)
      \item Nomes como \texttt{x}, \texttt{aux}, \texttt{dados2}
      \item Muitos níveis de \texttt{if}/\texttt{for} aninhados
    \end{itemize}
    \column{0.5\textwidth}
    \begin{itemize}
      \item Números ``mágicos'' soltos no código
      \item Variáveis globais modificadas em várias funções
      \item Comentários explicando código confuso
      \item Funções que leem do teclado \emph{e} calculam
    \end{itemize}
  \end{columns}
\end{frame}

\begin{frame}[fragile]{Antes e depois: nomes e números mágicos}
  \begin{columns}[T]
    \column{0.48\textwidth}
    \textbf{Antes}
\begin{lstlisting}[basicstyle=\ttfamily\scriptsize]
def f(d):
    t = 0
    for k in d:
        t = t + d[k][1] * d[k][2]
    if t > 5000:
        t = t * 0.95
    return t
\end{lstlisting}
    \column{0.5\textwidth}
    \textbf{Depois}
\begin{lstlisting}[basicstyle=\ttfamily\scriptsize]
LIMITE_DESCONTO = 5000.0
FATOR_DESCONTO = 0.95

def custo_total(estoque):
    """Custo do estoque, com desconto
    para compras acima do limite."""
    total = sum(item["preco"] * item["qtd"]
                for item in estoque.values())
    if total > LIMITE_DESCONTO:
        total *= FATOR_DESCONTO
    return total
\end{lstlisting}
  \end{columns}
  \vspace{1ex}
  \small Nomes que explicam, constantes nomeadas (Aula 5), docstring, compreensão no lugar de laço manual.
\end{frame}

\begin{frame}[fragile]{Antes e depois: eliminar repetição}
  \begin{columns}[T]
    \column{0.48\textwidth}
    \textbf{Antes}
\begin{lstlisting}[basicstyle=\ttfamily\scriptsize]
p = float(input("Preço: "))
while p < 0:
    p = float(input("Preço: "))
q = float(input("Quantidade: "))
while q < 0:
    q = float(input("Quantidade: "))
m = float(input("Mínimo: "))
while m < 0:
    m = float(input("Mínimo: "))
\end{lstlisting}
    \column{0.5\textwidth}
    \textbf{Depois}
\begin{lstlisting}[basicstyle=\ttfamily\scriptsize]
def le_nao_negativo(msg):
    while True:
        try:
            v = float(input(msg))
            if v >= 0:
                return v
        except ValueError:
            pass
        print("Valor inválido.")

preco = le_nao_negativo("Preço: ")
qtd = le_nao_negativo("Quantidade: ")
minimo = le_nao_negativo("Mínimo: ")
\end{lstlisting}
  \end{columns}
  \small \textbf{DRY} (\emph{don't repeat yourself}): se você copiou e colou, provavelmente falta uma função.
\end{frame}

\begin{frame}{Separação de responsabilidades}
  \centering
  \begin{tikzpicture}[caixa/.style={draw=uerjazul,fill=uerjcinza,rounded corners,minimum width=3.2cm,minimum height=1cm,align=center,font=\small},
      seta/.style={-{Stealth},thick,uerjazul}]
    \node[caixa] (io) at (0,0) {Interface\\\scriptsize menu, \texttt{input}, \texttt{print}};
    \node[caixa] (reg) at (4.5,0) {Regras\\\scriptsize cálculos, validações};
    \node[caixa] (dad) at (9,0) {Dados\\\scriptsize arquivos, CSV};
    \draw[seta,<->] (io)--(reg);
    \draw[seta,<->] (reg)--(dad);
  \end{tikzpicture}
  \vspace{1ex}

  \begin{itemize}\small
    \item \textbf{Regras} não deveriam usar \texttt{input}/\texttt{print}: recebem parâmetros e devolvem valores
    \item \textbf{Interface} não deveria calcular: chama as regras
    \item \textbf{Dados} sabem ler e gravar; não sabem o que significam os números
    \item Resultado: cada parte é testável e substituível
  \end{itemize}
\end{frame}

\begin{frame}[fragile]{Testes automáticos com \texttt{assert}}
\begin{lstlisting}[basicstyle=\ttfamily\scriptsize]
# testes.py
from estoque import custo_total, baixa_estoque, insumos_em_alerta

def novo():
    return {"cimento": {"unidade": "saco", "preco": 10.0, "qtd": 5, "minimo": 3}}

def teste_custo_total():
    assert custo_total(novo()) == 50.0
    assert custo_total({}) == 0.0                 # caso-limite: vazio

def teste_baixa():
    e = novo()
    baixa_estoque(e, "cimento", 2)
    assert e["cimento"]["qtd"] == 3
    assert insumos_em_alerta(e) == ["cimento"]    # qtd == minimo
\end{lstlisting}
\end{frame}

\begin{frame}[fragile]{Executando os testes}
\begin{lstlisting}[basicstyle=\ttfamily\scriptsize]
def teste_baixa_insuficiente():
    try:
        baixa_estoque(novo(), "cimento", 99)
        assert False, "deveria levantar ValueError"
    except ValueError:
        pass

if __name__ == "__main__":
    teste_custo_total()
    teste_baixa()
    teste_baixa_insuficiente()
    print("Todos os testes passaram!")
\end{lstlisting}
  \begin{itemize}\small
    \item Rode \texttt{python testes.py} após \alert{cada} alteração
    \item Para projetos maiores, existem \texttt{unittest} (biblioteca padrão) e \texttt{pytest} (\texttt{pip install pytest}): descobrem e executam todas as funções \texttt{test\_*} e mostram relatório
    \item Testes escritos viram a \textbf{rede de segurança} das refatorações
  \end{itemize}
\end{frame}

\begin{frame}{Casos extremos: o que testar?}
  \centering\small
  \begin{tabular}{p{3.4cm}p{8.4cm}}
    \toprule
    \textbf{Tipo} & \textbf{Exemplos}\\
    \midrule
    Vazio & lista/dicionário/arquivo vazio; texto sem nada\\
    Limites & zero, um, o menor e o maior valor permitido, igual ao mínimo\\
    Inválidos & texto onde se espera número; número negativo; opção inexistente\\
    Tamanho & um só elemento; milhares de linhas\\
    Repetição & item duplicado; cadastrar duas vezes o mesmo nome\\
    Arquivos & inexistente; sem permissão; linha incompleta; acentos\\
    Numéricos & divisão por zero; arredondamento de \texttt{float}\\
    \bottomrule
  \end{tabular}
\end{frame}

\begin{frame}{Ferramentas de qualidade}
  \begin{columns}[T]
    \column{0.5\textwidth}
    \begin{block}{Estilo (Aula 24)}
      \texttt{pylint programa.py}\\
      \texttt{yapf -i programa.py}\\
      \texttt{black programa.py}
    \end{block}
    \column{0.5\textwidth}
    \begin{block}{Revisão}
      \begin{itemize}\small
        \item Leia o código do colega \emph{como um estranho}
        \item Pergunte: dá para entender sem explicação?
        \item Troque de lugar: um integrante tenta rodar o programa do outro
      \end{itemize}
    \end{block}
  \end{columns}
  \vspace{1ex}
  \begin{alertblock}{Código ruim se apaga}
    Se um trecho está cada vez mais remendado, reescreva-o do zero (Cap.~9.4), agora que você conhece o problema.
  \end{alertblock}
\end{frame}

\begin{frame}{Atividade: sessão de refatoração e testes}
  \begin{enumerate}
    \item Escrevam \textbf{pelo menos 5 testes automáticos} das funções de regra (inclua 2 casos-limite)
    \item Passem o \texttt{pylint} (ou \texttt{yapf}) no código e tratem os principais avisos
    \item Procurem e eliminem \textbf{código repetido}
    \item Renomeiem variáveis e funções pouco claras; acrescentem \textbf{docstrings}
    \item Verifiquem que nenhuma função de regra usa \texttt{input}/\texttt{print}
    \item Mostrem ao professor: o que mudou? Os testes continuam passando?
  \end{enumerate}
\end{frame}

\begin{frame}{Checklist final antes da apresentação}
  \small
  \begin{tabular}{@{}cl@{}}
    $\square$ & O programa roda do zero com um único comando (e há instruções de uso)?\\
    $\square$ & Existem testes e todos passam?\\
    $\square$ & Nenhuma função tem mais de ~30 linhas ou faz várias coisas?\\
    $\square$ & Não há código repetido, números mágicos nem variáveis globais mutáveis?\\
    $\square$ & Entradas inválidas e erros de arquivo não derrubam o programa?\\
    $\square$ & Nomes, docstrings e comentários estão claros?\\
    $\square$ & Os dados de exemplo e os arquivos necessários estão na pasta do projeto?\\
    $\square$ & Cada integrante consegue explicar qualquer parte do código?\\
  \end{tabular}
\end{frame}

\end{document}
